\chapter{Рабочий пакет ветви \texorpdfstring{\(B-L\)}{B-L}-root и неоднородного дефекта}

Настоящий раздел сводит результаты отдельных гейтов в одну статусную
формулу. Подтверждение относится к явно проверенному минимальному меню и не
является доказательством абсолютной единственности среди всех возможных
расширений.

\section{Точно подтверждённые узлы}

В левокиральной вейлевской конвенции добавление одного стерильного поля
\(N^c\) на поколение с зарядом \(B-L=1\) зануляет все шесть локальных
непрерывных аномалий:
\begin{equation}
 \mathcal A_{(B-L)^3}=\mathcal A_{\mathrm{grav}^2(B-L)}
 =\mathcal A_{SU(3)^2(B-L)}=\mathcal A_{SU(2)^2(B-L)}
 =\mathcal A_{Y^2(B-L)}=\mathcal A_{Y(B-L)^2}=0.
\end{equation}
При
\begin{equation}
 \oint_y a_{B-L}=\frac{\pi}{2}
\end{equation}
стерильная линия получает корневую голономию
\begin{equation}
 \operatorname{Hol}_{N^c}(y)=i.
\end{equation}
Поле спаривания \(\Phi_{B-L}\) с зарядом \(-2\) допускает
калибровочно-инвариантную вершину
\begin{equation}
 y_N\Phi_{B-L}N^cN^c,
 \qquad -2+1+1=0.
\end{equation}

На меридиане корневой поток и ненулевой конденсат дают точное условное
следствие
\begin{equation}
 \oint_{\mu_\gamma}a_{B-L}=\pi,\quad |\Phi_{B-L}|\neq0
 \quad\Longrightarrow\quad
 \left|\operatorname{wind}(\Phi_{B-L})\right|=1.
\end{equation}
Слово ``условное'' существенно: топология фиксирует намотку при наличии
конденсата, но не вынуждает сам конденсат и не выбирает знак после смены
ориентации меридиана.

Для минимального функционала Гинзбурга--Ландау ковариантные импульсы равны
\begin{equation}
 k_n=\frac{2\pi n+\pi}{L}.
\end{equation}
В проверенном седловом анзаце с постоянным модулем ненулевая стационарная
конфигурация существует тогда и только тогда, когда
\begin{equation}
 \lambda v^2>\frac{\pi^2}{L^2}.
\end{equation}
При \(L=\pi\) это даёт точный порог \(\lambda v^2>1\). Нижние ветви
\(n=0\) и \(n=-1\) имеют импульсы \(k=+1\) и \(k=-1\), одинаковую энергию
и образуют точно вырожденную сопряжённую пару.

\begin{proposition}[Выжившая конструктивная ветвь минимального меню]
После проведённых гейтов единственной не закрытой конструктивной ветвью
проверенного минимального меню является
\begin{equation}
 U(1)_{B-L}+N^c+\Phi_{-2}
 \quad\longrightarrow\quad
 \text{неоднородная дефектная текстура при }\lambda v^2>1.
\end{equation}
Она подтверждена на уровне представлений, локальных аномалий, корневой
голономии и условного стационарного решения, но ещё не выведена из единого
спектрального действия.
\end{proposition}

\section{Уточнение роли торсионного скручивания}

Торсионное скручивание не является полным спасением однородной модели.
Оно меняет голономию поля спаривания с \(-1\) на \(+1\), но тогда линейная
вершина Юкавы получает остаточный торсионный знак \(-1\). Поэтому полностью
скрученный однородный конденсат не удовлетворяет одновременно условиям
корневой голономии, скалярной вершины и ненулевого параллельного вакуума.
Выжившую конфигурацию корректно называть неоднородной \(B-L\)-дефектной
текстурой. Термин ``торсионно-скрученная'' допустим только после вывода
соответствующего торсионно-нечётного отображения Юкавы, которого в текущем
действии нет.

\section{Две независимые точные конструкции}

Шестиканальная гауссова модель с обычной суммой по компонентам даёт точную
идентичность
\begin{equation}
 \operatorname{Var}(\bar x)
 =\frac{1}{6\zeta(4,1/2)}
 =\frac{1}{\pi^4}.
\end{equation}
Независимо от неё совместная \(\mathbb Z_2/\mathbb Z_4\)-проекция
вектороподобного \(SU(5)\)-родителя оставляет нулевые моды
\(U+2D+H\) и точно даёт направление
\begin{equation}
 \Delta b_{Y,2,3}=\left(\frac{17}{6},\frac16,2\right).
\end{equation}
Обе формулы воспроизводимы, но пока не следуют из одного оператора, одной
меры или одного действия. Поэтому запрещено записывать их композицию как
единый физический вывод. Первая конструкция остаётся кандидатом на
обратную восприимчивость, вторая --- точной проекцией представлений с
невыведенной пороговой величиной и неверным знаком обычного промежуточного
ренормгруппового пробега.

\section{Невыведенные динамические данные}

Топология и голономия не фиксируют:
\begin{itemize}
 \item величину \(\lambda v^2\) и сам факт превышения порога;
 \item выбор одной из ориентаций \(n=0,-1\);
 \item масштаб нарушения \(v_{B-L}\) и матрицу \(y_N\);
 \item граничное значение и пробег кинетического смешивания двух
 абелевых факторов.
\end{itemize}
Последний пункт особенно важен, поскольку
\begin{equation}
 \sum_{\mathrm{Weyl}}Y(B-L)=\frac83
\end{equation}
на поколение, и смешивание в общем случае радиативно порождается, но его
нормировка не определяется топологией.

\section{Закрытые маршруты}

Без новой динамики повторно не открываются следующие гипотезы:
\begin{enumerate}
 \item разложение \(1/3=1/4+1/12\) как вывод вещественного массового
 сдвига \(-\alpha/3\) из APS-фазы;
 \item универсальное значение \(\eta_D=-1/2\);
 \item отождествление \(|\zeta_R(-1)|=1/12\) с окружным
 логарифмическим определителем;
 \item получение \(23=24-1\) удалением калибровочной или конформной моды;
 \item использование трёхформы Черна--Саймонса
 \(\operatorname{Tr}(A\wedge dA+\tfrac23A^3)\) как пятимерного объёмного
 действия;
 \item получение стерильного корня из минимального \(SU(5)\), спиновой
 структуры или намбу-удвоения;
 \item однородный конденсат обычного поля заряда \(-2\);
 \item торсионное скручивание как одновременное спасение конденсата и
 вершины Юкавы;
 \item электрослабое и КХД-замыкание через обычный split-running или
 минимальный логарифмический конус порогов;
 \item формула \(23=k+1\) с \(k=22\) как выведенный уровень
 Черна--Саймонса или Весса--Зумино--Виттена.
\end{enumerate}

\begin{protocol}[Условие следующего продвижения]
Ветвь повышает статус только после вывода из одного заранее фиксированного
действия величины \(\lambda v^2\), ориентационно-нечётного члена,
\(v_{B-L}\), нормировки кинетического смешивания и оператора
Боголюбова--де Жена на выбранной конфигурации. Новые численные совпадения
без такого вывода не считаются продолжением ветви.
\end{protocol}